\chapter{相变与临界现象基础}

\section{相平衡条件}

两相共存时，热平衡、力学平衡、化学平衡分别要求
\[
  T_1=T_2,\qquad P_1=P_2,\qquad \mu_1=\mu_2.
\]
因此单组分系统的相图可由化学势比较决定：稳定相是给定 $T,P$ 下 Gibbs 自由能最低的相。

\section{Clapeyron 方程}

沿相平衡曲线有
\[
  \mu_1(T,P)=\mu_2(T,P).
\]
利用
\[
  \dd\mu=-s\dd T+v\dd P
\]
得到
\[
  \frac{\dd P}{\dd T}=\frac{s_2-s_1}{v_2-v_1}
  =\frac{L}{T(v_2-v_1)}.
\]
其中 $L=T(s_2-s_1)$ 为单位粒子的潜热。

\section{一级相变与连续相变}

一级相变的特征是熵或体积等一阶导数不连续，通常有潜热。连续相变没有潜热，但热容、压缩率、磁化率等二阶导数可发散或不连续。

Landau 理论用序参量 $\eta$ 描述对称性变化：
\[
  F(T,\eta)=F_0(T)+a(T)\eta^2+b\eta^4+\cdots,\qquad b>0.
\]
若 $a(T)=a_0(T-T_c)$，则
\[
  \eta=0\quad (T>T_c),\qquad
  \eta^2=-\frac{a(T)}{2b}\quad (T<T_c).
\]

\begin{keybox}
\textbf{考试抓手：}相变题通常先写共存条件 $\mu_1=\mu_2$，再对两边微分推出斜率关系。不要从记忆中的水蒸气公式直接硬套。
\end{keybox}

